\documentclass{article}
\usepackage{amsmath,amssymb}

\begin{document}

% Extracted from Q34_Version3.tex

% item=1 kind=display context=plain lines=97-99
\[
\displaystyle\int \frac{x^2}{(4-x^2)^{\frac{3}{2}}}\, dx.
\]

% item=2 kind=display context=solutionStep:body lines=125-127
\[
\frac{x^2}{(4-x^2)^\frac{3}{2}},
\]

% item=3 kind=display context=solutionStep:body lines=129-131
\[
(4-x^2)^{\frac{3}{2}}.
\]

% item=4 kind=display context=solutionStep:body lines=137-139
\[
(4-x^2)^{\frac{3}{2}}
\]

% item=5 kind=display context=solutionStep:body lines=141-143
\[
a^{\frac{m}{n}}=\sqrt[n]{a^m}=\left(\sqrt[n]{a}\right)^m,\qquad \sqrt{a} =a^{\frac{1}{2}}.
\]

% item=6 kind=display context=solutionStep:body lines=148-168
\[
\begin{array}{@{}r@{}l@{}}
 Rewrite \frac{3}{2}=\frac{1}{2}\cdot 3
 &
 (4-x^2)^{\frac{3}{2}}
 = (4-x^2)^{\frac{1}{2}\cdot 3}
 \\[2.2ex]
 Use a^{bc}=(a^b)^c
 &
 (4-x^2)^{\frac{1}{2}\cdot 3}
 =
 \left((4-x^2)^{\frac{1}{2}}\right)^3
 \\[2.2ex]
 Rewrite (4-x^2)^{\frac{1}{2}} as a square root &
 \left((4-x^2)^{\frac{1}{2}}\right)^3
 =
 \left(\sqrt{4-x^2}\right)^3
 \end{array}
\]

% item=7 kind=display context=solutionStep:body lines=183-185
\[
\sqrt{4-x^2},
\]

% item=8 kind=display context=solutionStep:body lines=187-189
\[
\sqrt{a^2-x^2}.
\]

% item=9 kind=display context=solutionStep:body lines=200-202
\[
\sqrt{a^2-x^2},
\]

% item=10 kind=display context=solutionStep:body lines=204-206
\[
x=a\sin\theta.
\]

% item=11 kind=display context=solutionStep:body lines=211-219
\[
\begin{aligned}
 Substitute x=a\sin\theta
 \qquad &
 \sqrt{a^2-x^2}
 =
 \sqrt{a^2-(a\sin\theta)^2}
 \end{aligned}
\]

% item=12 kind=display context=solutionStep:body lines=227-242
\[
\begin{array}{@{}r@{}l@{}}
 Use Power of a Product Rule , &
 (a b )^2=a^2 b ^2
 \\[1.2ex]
 Substitute b=\sin\theta,
 &
 (a \sin\theta )^2=a^2( \sin\theta )^2
 \\[1.2ex]
 &
 \sqrt{a^2- (a\sin\theta)^2 }
 = \sqrt{a^2- a^2(\sin\theta)^2 }
 \end{array}
\]

% item=13 kind=environment:aligned context=solutionStep:body lines=247-250
\[
\begin{aligned}[t]
 For pattern \sqrt{a^2-x^2}, let x &= a\sin\theta\\
 \sqrt{a^2-x^2} &= \sqrt{a^2-a^2(\sin\theta)^2}
 \end{aligned}
\]

% item=14 kind=display context=solutionStep:body lines=258-295
\[
\begin{array}{@{}r@{}l@{}r@{}c@{}l@{}}
 Rewrite 4=2\cdot 2=
 &
 2^2
 &
 \sqrt{4-x^2}
 &
 =
 &
 \sqrt{2^2-x^2}
 \\[1.2ex]
 &
 Let &
 \sqrt{2^2-x^2}
 &
 =
 &
 \sqrt{a^2-x^2}
 \\[1.2ex]
 &
 &
 a
 &
 =
 &
 2 \\[1.2ex]
 &
 &
 x
 &
 =
 &
 a \sin\theta = 2 \sin\theta
 \end{array}
\]

% item=15 kind=environment:aligned context=solutionStep:body lines=303-308
\[
\begin{aligned}[t]
 \sqrt{a^2-x^2}
 &= \sqrt{a^2-a^2(\sin\theta)^2}\\
 \sqrt{2^2-x^2}
 &= \sqrt{2^2-2^2(\sin\theta)^2}
 \end{aligned}
\]

% item=16 kind=display context=solutionStep:body lines=316-323
\[
\begin{array}{@{}r@{}l@{}}
 Replace 2^2=4
 &
 \sqrt{ 2^2 - 2^2 (\sin\theta)^2}
 = \sqrt{ 4 - 4 (\sin\theta)^2}
 \end{array}
\]

% item=17 kind=display context=solutionStep:body lines=328-340
\[
\begin{array}{@{}r@{}l@{}}
 Rewrite 4=4\cdot 1
 &
 4 -4(\sin\theta)^2
 = 4(1) -4((\sin\theta)^2)
 \\[1.2ex]
 Factor out 4
 &
 4 - 4 (\sin\theta)^2
 = 4 (1-(\sin\theta)^2)
 \end{array}
\]

% item=18 kind=display context=solutionStep:body lines=345-372
\[
\begin{array}{@{}r@{}r@{}c@{}l@{}}
 By (\sin\theta)^2=\sin^2\theta , &
 4(1- (\sin\theta)^2 )
 &
 =
 &
 4(1- \sin^2\theta )
 \\[1.2ex]
 By Pythagorean identity , &
 1-\sin^2\theta
 &
 =
 &
 \cos^2\theta
 \\[1.2ex]
 Substitute the Pythagorean identity, &
 4( 1-\sin^2\theta )
 &
 =
 &
 4( \cos^2\theta )
 \end{array}
\]

% item=19 kind=environment:aligned context=solutionStep:body lines=377-383
\[
\begin{aligned}[t]
 \sqrt{4-x^2}
 &= \sqrt{2^2-x^2}\quad from Step 3.1 \\
 &= \sqrt{2^2-2^2(\sin\theta)^2}\quad from Step 3.2 \\
 &= \sqrt{4(1-\sin^2\theta)}\quad from Step 4.1 and Step 4.2 \\
 &= \sqrt{4\cos^2\theta}\quad from Step 4.3 \end{aligned}
\]

% item=20 kind=display context=solutionStep:body lines=391-404
\[
\begin{array}{@{}r@{}l@{}}
 From Step 1.3 , &
 (4-x^2)^{\frac{3}{2}}=\left(\sqrt{4-x^2}\right)^3
 \\[1.5ex]
 Substitute \sqrt{4-x^2}=\sqrt{4\cos^2\theta}
 &
 \left(\sqrt{4-x^2}\right)^3
 =
 \left(\sqrt{4\cos^2\theta}\right)^3
 \end{array}
\]

% item=21 kind=display context=solutionStep:body lines=412-426
\[
\begin{array}{@{}r@{}l@{}}
 Since \sqrt{ab}=\sqrt{a}\sqrt{b} , &
 \sqrt{4\,\cos^2\theta}
 =
 \sqrt4\,\sqrt{\cos^2\theta}
 \\[2.0ex]
 \sqrt4=2 &
 \sqrt4\,\sqrt{\cos^2\theta}
 =
 2\,\sqrt{\cos^2\theta}
 \end{array}
\]

% item=22 kind=display context=solutionStep:body lines=435-447
\[
\begin{array}{@{}r@{}l@{}}
 Use \sqrt{u^2}=|u|
 &
 \sqrt{\cos^2\theta}=|\cos\theta|
 \\[1.2ex]
 Multiply both sides by 2
 &
 2 \sqrt{\cos^2\theta}
 =
 2 |\cos\theta|
 \end{array}
\]

% item=23 kind=display context=solutionStep:body lines=453-473
\[
\begin{array}{@{}r@{}l@{}}
 Cube both sides &
 \left(\sqrt{4\cos^2\theta}\right)^{ 3 }
 =
 (2|\cos\theta|)^{ 3 }
 \\[1.2ex]
 Use (ab)^3=a^3b^3
 &
 (2|\cos\theta|)^3
 =
 2^3 |\cos\theta|^3
 \\[1.2ex]
 Simplify 2^3
 &
 \left(\sqrt{4\cos^2\theta}\right)^3
 =
 8 |\cos\theta|^3
 \end{array}
\]

% item=24 kind=display context=solutionStep:body lines=478-489
\[
\begin{array}{@{}r@{}l@{}}
 For x=2\sin\theta , we take 0\le \theta\le \frac{\pi}{2}
 &
 \cos\theta\ge 0
 \\[1.2ex]
 So |\cos\theta|=\cos\theta
 &
 8|\cos\theta|^3=8\cos^3\theta
 \end{array}
\]

% item=25 kind=display context=solutionStep:body lines=493-504
\[
\begin{aligned}[t]
 \sqrt{4\cos^2\theta}
 &= 2\sqrt{\cos^2\theta}\quad from Step 5.2 \\
 2\sqrt{\cos^2\theta}
 &= 2|\cos\theta|\quad from Step 5.3 \\
 \left(\sqrt{4\cos^2\theta}\right)^3
 &= 8|\cos\theta|^3\quad from Step 5.4 \\
 (4-x^2)^{\frac{3}{2}}
 &= 8\cos^3\theta\quad from Step 5.6 \end{aligned}
\]

% item=26 kind=display context=solutionStep:body lines=512-522
\[
\begin{array}{@{}r@{}l@{}}
 Start with the substitution &
 x=2\sin\theta\quad from Step 3.1 \\[1.2ex]
 Square both sides &
 x^2=(2\sin\theta)^2
 \end{array}
\]

% item=27 kind=display context=solutionStep:body lines=527-539
\[
\begin{array}{@{}r@{}l@{}}
 Let a=2,\; b=\sin\theta
 &
 (ab)^2=a^2b^2
 \\[1.2ex]
 Substitute into the rule &
 (2\sin\theta)^{ 2 }
 =
 2^{ 2 }(\sin\theta)^{ 2 }
 \end{array}
\]

% item=28 kind=display context=solutionStep:body lines=545-563
\[
\begin{array}{@{}r@{}l@{}}
 Compute 2^2
 &
 2^2=4
 \\[1.2ex]
 By square of trigonometric functions &
 (\sin\theta)^2=\sin^2\theta
 \\[1.2ex]
 Substitute &
 2^2(\sin\theta)^2=4\sin^2\theta
 \\[1.2ex]
 Therefore &
 x^2=4\sin^2\theta
 \end{array}
\]

% item=29 kind=display context=solutionStep:body lines=571-594
\[
\begin{array}{@{}r@{}l@{}}
 Start with the substitution &
 x=2\sin\theta\quad from Step 3.1 \\[1.2ex]
 Differentiate both sides with respect to \theta
 &
 \frac{dx}{d\theta}=\frac{d}{d\theta}(2\sin\theta)
 \\[1.2ex]
 Use the constant multiple rule &
 \frac{dx}{d\theta}
 =
 2 \frac{d}{d\theta}(\sin\theta)
 \\[1.2ex]
 Since \frac{d}{d\theta}(\sin\theta)=\cos\theta
 &
 \frac{dx}{d\theta}
 =
 2 \cos\theta
 \end{array}
\]

% item=30 kind=display context=solutionStep:body lines=600-616
\[
\begin{array}{@{}r@{}l@{}}
 Start with &
 \frac{dx}{d\theta}=2\cos\theta
 \\[1.2ex]
 Multiply both sides by d\theta
 &
 d\theta \left(\frac{dx}{d\theta}\right)
 =
 (2\cos\theta) d\theta \\[1.2ex]
 So &
 dx=2\cos\theta\, d\theta
 \end{array}
\]

% item=31 kind=display context=solutionStep:body lines=624-630
\[
\begin{aligned}[t]
 x^2 &= 4\sin^2\theta \quad from Step 6 \\
 (4-x^2)^{\frac{3}{2}} &= 8\cos^3\theta \quad from Step 1 to Step 5 \\
 dx &= 2\cos\theta\, d\theta \quad from Step 7 \end{aligned}
\]

% item=32 kind=display context=solutionStep:body lines=635-650
\[
\begin{array}{@{}r@{}l@{}}
 Start with the original integral &
 \displaystyle\int \frac{x^2}{(4-x^2)^{\frac{3}{2}}}\,dx
 \\[1.2ex]
 Substitute the expressions from Step 8.1 &
 \displaystyle\int \frac{x^2}{(4-x^2)^{\frac{3}{2}}}\,dx
 =
 \displaystyle\int
 \frac{ 4\sin^2\theta }{ 8\cos^3\theta }\,
 (2\cos\theta\, d\theta) \end{array}
\]

% item=33 kind=display context=solutionStep:body lines=658-673
\[
\begin{array}{@{}r@{}l@{}}
 &
 \displaystyle\int \frac{4\sin^2\theta}{8\cos^3\theta}\, 2\cos\theta \, d\theta
 =
 \displaystyle\int \frac{4\cdot 2 \sin^2\theta \cos\theta }{8\cos^3\theta}\, d\theta
 \\[1.2ex]
 Simplify 4\cdot 2
 &
 \displaystyle\int \frac{4\cdot2\sin^2\theta\cos\theta}{8\cos^3\theta}\, d\theta
 =
 \displaystyle\int \frac{ 8 \sin^2\theta\cos\theta}{8\cos^3\theta}\, d\theta
 \end{array}
\]

% item=34 kind=display context=solutionStep:body lines=678-687
\[
\begin{array}{@{}r@{}l@{}}
 Cancel the common factor 8
 &
 \displaystyle\int \frac{ 8 \sin^2\theta\cos\theta}{ 8 \cos^3\theta}\, d\theta
 =
 \displaystyle\int \frac{\sin^2\theta\cos\theta}{\cos^3\theta}\, d\theta
 \end{array}
\]

% item=35 kind=display context=solutionStep:body lines=692-711
\[
\begin{array}{@{}r@{}l@{}}
 Rewrite \cos^3\theta
 &
 \cos^3\theta=\cos^2\theta\cos\theta
 \\[1.2ex]
 Substitute into the denominator &
 \displaystyle\int \frac{\sin^2\theta\cos\theta}{ \cos^3\theta }\, d\theta
 =
 \displaystyle\int \frac{\sin^2\theta\cos\theta}{ \cos^2\theta\cos\theta }\, d\theta
 \\[1.2ex]
 Cancel \cos\theta
 &
 \displaystyle\int \frac{\sin^2\theta \cos\theta }{\cos^2\theta \cos\theta }\, d\theta
 =
 \displaystyle\int \frac{\sin^2\theta}{\cos^2\theta}\, d\theta
 \end{array}
\]

% item=36 kind=display context=solutionStep:body lines=716-732
\[
\begin{array}{@{}r@{}l@{}}
 By definition of tangent, &
 \tan\theta=\dfrac{\sin\theta}{\cos\theta}
 \\[1.2ex]
 Square both sides &
 \tan^2\theta=\dfrac{\sin^2\theta}{\cos^2\theta}
 \\[1.2ex]
 Rewrite the integrand &
 \displaystyle\int \dfrac{\sin^2\theta}{\cos^2\theta}\, d\theta
 =
 \displaystyle\int \tan^2\theta\, d\theta
 \end{array}
\]

% item=37 kind=display context=solutionStep:body lines=740-759
\[
\begin{array}{@{}r@{}l@{}}
 Use the trigonometry identity &
 \tan^2\theta=\sec^2\theta-1
 \\[1.2ex]
 Substitute into the integral &
 \displaystyle\int \tan^2\theta \, d\theta
 =
 \displaystyle\int (\sec^2\theta-1) \, d\theta
 \\[1.2ex]
 By additive property of integrals &
 \displaystyle\int (\sec^2\theta-1)\, d\theta
 =
 \displaystyle\int \sec^2\theta\, d\theta-\displaystyle\int 1\, d\theta
 \end{array}
\]

% item=38 kind=display context=solutionStep:body lines=764-785
\[
\begin{array}{@{}r@{}l@{}}
 Integrate \sec^2\theta
 &
 \displaystyle\int \sec^2\theta\, d\theta=\tan\theta
 \\[1.2ex]
 Integrate 1
 &
 \displaystyle\int 1\, d\theta=\theta
 \\[1.2ex]
 Combine the result &
 \displaystyle\int \sec^2\theta\, d\theta-\displaystyle\int 1\, d\theta
 =
 \tan\theta-\theta+C
 \\[1.2ex]
 Therefore &
 \displaystyle\int \tan^2\theta\, d\theta=\tan\theta-\theta+C
 \end{array}
\]

% item=39 kind=display context=solutionStep:body lines=799-820
\[
\begin{array}{@{}r@{}l@{}}
 Start with &
 x=2\sin\theta\quad from Step 3.1 \\[1.2ex]
 Multiply both sides by \frac{1}{2}
 &
 \frac{1}{2} x
 =
 \frac{1}{2} 2\sin\theta
 \\[1.2ex]
 \frac{1}{2}\cdot 2=1 &
 \frac{x}{2}=\sin\theta
 \\[1.2ex]
 Apply \sin^{-1} to both side of the equation &
 \theta=\sin^{-1}\!\left(\frac{x}{2}\right)
 \end{array}
\]

% item=40 kind=display context=solutionStep:body lines=825-832
\[
\begin{array}{@{}r@{}l@{}}
 Use \sin\theta=\frac{ opposite }{ hypotenuse } &
 \sin\theta=\frac{ opposite }{ hypotenuse }=\frac{x}{2}
 \end{array}
\]

% item=41 kind=display context=solutionStep:body lines=837-844
\[
\begin{array}{@{}r@{}l@{}}
 From \sin\theta=\frac{x}{2}
 &
 opposite = x ,\quad hypotenuse = 2 \end{array}
\]

% item=42 kind=display context=solutionStep:body lines=859-870
\[
\begin{array}{@{}r@{}l@{}}
 By the Pythagorean Theorem , &
 adjacent =\sqrt{ hypotenuse ^2- opposite ^2}
 \\[1.2ex]
 Substitute the values &
 adjacent =\sqrt{2^2-x^2}= \sqrt{4-x^2} \end{array}
\]

% item=43 kind=display context=solutionStep:body lines=875-882
\[
\begin{array}{@{}r@{}l@{}}
 Use \tan\theta=\frac{ opposite }{ adjacent }
 &
 \tan\theta=\frac{ x }{ \sqrt{4-x^2} }
 \end{array}
\]

% item=44 kind=display context=solutionStep:body lines=890-908
\[
\begin{array}{@{}r@{}l@{}}
 From Step 10.2 , &
 \displaystyle\int \frac{x^2}{(4-x^2)^{\frac{3}{2}}}\,dx
 =
 \tan\theta-\theta+C
 \\[1.2ex]
 Substitute Step 11.2 and Step 11.6 &
 \displaystyle\int \frac{x^2}{(4-x^2)^{\frac{3}{2}}}\,dx
 =
 \frac{x}{\sqrt{4-x^2}} -
 \sin^{-1}\!\left(\frac{x}{2}\right) + C
 \end{array}
\]

% item=45 kind=display context=solutionStep:body lines=913-915
\[
The correct option is b .
\]

% item=46 kind=display context=plain lines=931-933
\[
\frac{d}{dx}\left(\frac{x^3}{3}\right)=x^2,
\]

% item=47 kind=display context=plain lines=942-944
\[
\frac{x^2}{(4-x^2)^{3/2}}
\]

% item=48 kind=display context=plain lines=953-955
\[
\frac{x^2}{(4-x^2)^{3/2}}
\]

% item=49 kind=display context=plain lines=965-967
\[
3x+7,
\]

% item=50 kind=display context=plain lines=977-979
\[
3x+7,
\]

% item=51 kind=display context=plain lines=986-988
\[
(ab)^n = a^n b^n.
\]

% item=52 kind=display context=plain lines=991-993
\[
(3x)^2=3^2x^2=9x^2.
\]

% item=53 kind=display context=plain lines=1002-1004
\[
\sqrt{4-x^2}=\sqrt{2^2-x^2}
\]

% item=54 kind=display context=plain lines=1008-1012
\[
\sqrt{a^2-x^2}: x=a\sin\theta,\qquad
 \sqrt{a^2+x^2}: x=a\tan\theta,\qquad
 \sqrt{x^2-a^2}: x=a\sec\theta.
\]

% item=55 kind=display context=plain lines=1018-1020
\[
a^{m/n}=\sqrt[n]{a^m}=\left(a^{1/n}\right)^m.
\]

% item=56 kind=display context=plain lines=1023-1025
\[
(4-x^2)^{3/2}=\left((4-x^2)^{1/2}\right)^3=\left(\sqrt{4-x^2}\right)^3.
\]

% item=57 kind=display context=plain lines=1031-1035
\[
(\sin\theta)^2=\sin^2\theta,\qquad
 (\cos\theta)^2=\cos^2\theta,\qquad
 (\tan\theta)^2=\tan^2\theta.
\]

% item=58 kind=display context=plain lines=1041-1045
\[
\sin^2\theta+\cos^2\theta=1
 \quad\Longleftrightarrow\quad
 1-\sin^2\theta=\cos^2\theta.
\]

% item=59 kind=display context=plain lines=1048-1050
\[
4-4\sin^2\theta=4(1-\sin^2\theta)=4\cos^2\theta.
\]

% item=60 kind=display context=plain lines=1056-1058
\[
a^2+b^2=c^2.
\]

% item=61 kind=display context=plain lines=1062-1064
\[
adjacent ^2=2^2-x^2=4-x^2.
\]

% item=62 kind=display context=plain lines=1069-1073
\[
y^2=x
 \quad\Longrightarrow\quad
 y is a square root of x.
\]

% item=63 kind=display context=plain lines=1075-1079
\[
\sqrt{x}=y
 \quad\Longleftrightarrow\quad
 y^2=x and y\ge 0.
\]

% item=64 kind=display context=plain lines=1082-1084
\[
\sqrt{9}=3
\]

% item=65 kind=display context=plain lines=1086-1088
\[
3^2=9.
\]

% item=66 kind=display context=plain lines=1091-1093
\[
\sqrt{x^2}=|x|
\]

% item=67 kind=display context=plain lines=1094-1096
\[
\sqrt{a}\sqrt{b}=\sqrt{ab}
\]

% item=68 kind=display context=plain lines=1097-1099
\[
\sqrt{\frac{a}{b}}=\frac{\sqrt{a}}{\sqrt{b}}
\]

% item=69 kind=display context=plain lines=1100-1104
\[
\left(\sqrt{a}\right)^2=a
 \quad and \quad
 a^\frac{1}{2}=\sqrt{a}
\]

% item=70 kind=display context=plain lines=1110-1112
\[
\frac{d}{d\theta}\bigl(c\,f(\theta)\bigr)=c\,f'(\theta).
\]

% item=71 kind=display context=plain lines=1115-1117
\[
\frac{d}{d\theta}(2\sin\theta)=2\cos\theta.
\]

% item=72 kind=display context=plain lines=1123-1125
\[
\int (f(x)+g(x))\,dx=\int f(x)\,dx + \int g(x)\,dx
\]

% item=73 kind=display context=plain lines=1127-1129
\[
\int (f(x)-g(x))\,dx=\int f(x)\,dx - \int g(x)\,dx.
\]

% item=74 kind=display context=plain lines=1132-1136
\[
\int (\sec^2\theta-1)\,d\theta
 =
 \int \sec^2\theta\,d\theta-\int 1\,d\theta.
\]

% item=75 kind=display context=plain lines=1144-1146
\[
|3|=3,\qquad |-3|=3.
\]

% item=76 kind=display context=plain lines=1148-1150
\[
\sqrt{u^2}=|u|.
\]

% item=77 kind=display context=plain lines=1152-1154
\[
|\cos\theta|=\cos\theta.
\]

% item=78 kind=display context=plain lines=1160-1164
\[
\sin\theta=\frac{ opposite }{ hypotenuse },\qquad
 \cos\theta=\frac{ adjacent }{ hypotenuse },\qquad
 \tan\theta=\frac{ opposite }{ adjacent }.
\]

% item=79 kind=display context=plain lines=1169-1171
\[
\sqrt{2^2-x^2}=\sqrt{4-x^2},
\]

% item=80 kind=display context=plain lines=1173-1175
\[
\tan\theta=\frac{x}{\sqrt{4-x^2}}.
\]

% item=81 kind=display context=plain lines=1184-1186
\[
\theta=\sin^{-1}\!\left(\frac{x}{2}\right).
\]

\end{document}
